\documentclass[10pt]{article}
\usepackage[margin=0.6in]{geometry}
\usepackage{amsmath,amssymb}
\usepackage{multicol}
\usepackage{booktabs}
\usepackage{array}
\usepackage{enumitem}
\usepackage[dvipsnames]{xcolor}

\pagestyle{empty}

% Tableaux system
\RequirePackage[tableaux]{prooftrees}
\RequirePackage{pifont}
\forestset{line numbering, 
        check with={\textcolor{ForestGreen}{\text{\ding{52}}}},
        close with format={red}}

% Open Logic commands for tableaux
\usepackage{open-logic-minimal}

% Define the tableau environment
\newenvironment{oltableau}{\center\tableau{}}{\endtableau\endcenter}

% Define True/False symbols
\usepackage[bb=boondox]{mathalfa}
\DeclareMathAlphabet{\mathbx}{U}{BOONDOX-ds}{m}{n}
\def\True{\mathbb{T}}
\def\False{\mathbb{F}}

% Spacing around tableaux
\usepackage{etoolbox}
\AfterEndEnvironment{oltableau}{\vspace{9pt}}
\BeforeBeginEnvironment{oltableau}{\vspace{9pt}}

% Color definitions (if you need them later)
\definecolor{darkblue}{rgb}{0,0,0.8}
\definecolor{dodgerblue}{RGB}{30,144,255}

% Assumption code
\newcommand*{\DeclareDocumentMacro}[2]{\def#1{#2}}
\DeclareDocumentMacro \TAss {Assumption}

\newcommand{\ruleheader}[1]{\medskip\noindent\textbf{\large #1}\par\smallskip\hrule\smallskip}
\newcommand{\rulesubheader}[1]{\bigskip\noindent\textbf{#1}\par\smallskip}

\setlength{\parindent}{0pt}

\begin{document}

\ruleheader{Part 1: S5}

Use the simplified rules for \textbf{S5} to build tableaux for each of these claims. Say which of them are valid, i.e., are theorems of \textbf{S5}.

\rulesubheader{1. $\Box(\Box p \rightarrow (q \vee r)) \vee \Box(\Box q \rightarrow p)$ in \textbf{S5}}

\begin{oltableau}
    [\sFmla{\False}{1, \Box (\Box p \rightarrow  (q \lor  r)) \lor  \Box (\Box q \rightarrow  p)}, checked, just = \TAss
        [\sFmla{\False}{1, \Box (\Box p \rightarrow  (q \lor  r))}, checked, just = {\TRule{\False}{\lor}[1]}
            [\sFmla{\False}{1, \Box (\Box q \rightarrow  p)}, checked, just = {\TRule{\False}{\lor}[1]}
                [\sFmla{\False}{2, \Box p \rightarrow  (q \lor  r)}, checked, just = {\TRule{\False}{\Box}[2]}
                    [\sFmla{\True}{2, \Box p}, just = {\TRule{\False}{\rightarrow}[4]}
                        [\sFmla{\False}{2, q \lor  r}, checked, just = {\TRule{\False}{\rightarrow}[4]}
                            [\sFmla{\False}{2, q}, just = {\TRule{\False}{\lor}[6]}
                                [\sFmla{\False}{2, r}, just = {\TRule{\False}{\lor}[2]}
                                    [\sFmla{\False}{3, \Box q \rightarrow  p}, just = {\TRule{\False}{\Box}[3]}
                                        [\sFmla{\True}{3, \Box q}, just = {\TRule{\False}{\rightarrow}[9]}
                                            [\sFmla{\False}{3, p}, just = {\TRule{\False}{\rightarrow}[9]}
                                                [\sFmla{\True}{3, p}, just = {\TRule{\True}{\Box}[5]}, close]
                                            ]
                                        ]
                                    ]
                                ]
                            ]
                        ]
                    ]
                ]
            ]
        ]
    ]
\end{oltableau}

The only branch closes, so this is \textbf{valid} in \textbf{S5}.

\rulesubheader{$\Diamond \Box (p \rightarrow q) \rightarrow (\Box p \rightarrow \Diamond q)$ in \textbf{S5}}

\begin{oltableau}
    [\sFmla{\False}{1, \Diamond \Box (p \rightarrow  q) \rightarrow  (\Box p \rightarrow  \Diamond q)}, checked, just = \TAss
        [\sFmla{\True}{1, \Diamond \Box (p \rightarrow  q)}, checked, just = {\TRule{\False}{\rightarrow}[1]}
            [\sFmla{\False}{1, \Box p \rightarrow  \Diamond q}, checked, just = {\TRule{\False}{\rightarrow}[1]}
                [\sFmla{\True}{1, \Box p}, just = {\TRule{\False}{\rightarrow}[3]}
                    [\sFmla{\False}{1, \Diamond q}, just = {\TRule{\False}{\rightarrow}[3]}
                        [\sFmla{\True}{2, \Box (p \rightarrow  q)}, just = {\TRule{\True}{\Diamond}[2]}
                            [\sFmla{\True}{2, p \rightarrow  q}, checked, just = {\TRule{\True}{\Box}[2]}
                                [\sFmla{\True}{2, p}, just = {\TRule{\True}{\Box}[4]}
                                    [\sFmla{\False}{2, q}, just = {\TRule{\False}{\Diamond}[5]}
                                        [\sFmla{\False}{2, p}, just = {\TRule{\True}{\rightarrow}[7]}, close]
                                        [\sFmla{\True}{2, q}, just = {\TRule{\True}{\rightarrow}[7]}, close]
                                    ]
                                ]
                            ]
                        ]
                    ]
                ]
            ]
        ]
    ]
\end{oltableau}

Both branches clos, so this is \textbf{valid} in \textbf{S5}.

\newpage

\rulesubheader{$\Box (p \vee q) \rightarrow (\Diamond (\neg q \vee r) \rightarrow \Diamond (p \vee r))$ in \textbf{S5}}

\begin{oltableau}
    [\sFmla{\False}{1, \Box (p \lor  q) \rightarrow  (\Diamond (\neg q \lor  r) \rightarrow  \Diamond (p \lor  r))}, checked, just = \TAss
        [\sFmla{\True}{1, \Box (p \lor  q)}, just = {\TRule{\False}{\rightarrow}[1]}
            [\sFmla{\False}{1, \Diamond (\neg q \lor  r) \rightarrow  \Diamond (p \lor  r)}, checked, just = {\TRule{\False}{\rightarrow}[1]}
                [\sFmla{\True}{1, \Diamond (\neg q \lor  r)}, checked, just = {\TRule{\False}{\rightarrow}[3]}
                    [\sFmla{\False}{1, \Diamond (p \lor  r)}, just = {\TRule{\False}{\rightarrow}[3]}
                        [\sFmla{\True}{2, \neg q \lor  r}, checked, just = {\TRule{\True}{\Diamond}[4]}
                            [\sFmla{\True}{2, p \lor  q}, checked, just = {\TRule{\True}{\Box}[1]}
                                [\sFmla{\False}{2, p \lor  r}, checked, just = {\TRule{\False}{\Diamond}[5]}
                                    [\sFmla{\False}{2, p}, just = {\TRule{\False}{\lor}[8]}
                                        [\sFmla{\False}{2, r}, just = {\TRule{\False}{\lor}[8]}
                                            [\sFmla{\True}{2, p}, just = {\TRule{\True}{\lor}[7]}, close]
                                            [\sFmla{\True}{2, q}, just = {\TRule{\True}{\lor}[7]}
                                                [\sFmla{\True}{2, \neg q}, checked, just = {\TRule{\True}{\lor}[6]}
                                                    [\sFmla{\False}{2, q}, just = {\TRule{\True}{\neg}[12]}, close]
                                                ]
                                                [\sFmla{\True}{2, r}, just = {\TRule{\True}{\lor}[6]}, close]
                                            ]
                                        ]
                                    ]
                                ]
                            ]
                        ]
                    ]
                ]
            ]
        ]
    ]
\end{oltableau}

All branches close, so this is \textbf{valid} in \textbf{S5}.

\rulesubheader{$(\neg \Diamond p \vee \neg \Diamond \Diamond p) \rightarrow \neg \Diamond \Diamond \Diamond p$ in \textbf{S5}}

\begin{oltableau}
    [\sFmla{\False}{1, (\neg \Diamond p \lor  \neg \Diamond \Diamond p) \rightarrow  \neg \Diamond \Diamond \Diamond p}, checked, just = \TAss
        [\sFmla{\True}{1, \neg \Diamond p \lor  \neg \Diamond \Diamond p}, checked, just = {\TRule{\False}{\rightarrow}[1]}
            [\sFmla{\False}{1, \neg \Diamond \Diamond \Diamond p}, checked, just = {\TRule{\False}{\rightarrow}[1]}
                [\sFmla{\True}{1, \Diamond \Diamond \Diamond p}, checked, just = {\TRule{\False}{\neg}[3]}
                    [\sFmla{\True}{2, \Diamond \Diamond p}, checked, just = {\TRule{\True}{\Diamond}[4]}
                        [\sFmla{\True}{3, \Diamond p}, checked, just = {\TRule{\True}{\Diamond}[5]}
                            [\sFmla{\True}{4, p}, just = {\TRule{\True}{\Diamond}[6]}
                                [\sFmla{\True}{1, \neg \Diamond p}, checked, just = {\TRule{\True}{\lor}[2]}
                                    [\sFmla{\False}{1, \Diamond p}, just = {\TRule{\True}{\neg}[8]}
                                        [\sFmla{\False}{4, p}, just = {\TRule{\False}{\Diamond}[9]}, close]
                                    ]
                                ]
                                [\sFmla{\True}{1, \neg \Diamond \Diamond p}, checked, just = {\TRule{\True}{\lor}[2]}
                                    [\sFmla{\False}{1, \Diamond \Diamond p}, just = {\TRule{\True}{\neg}[8]}, move by = 2
                                        [\sFmla{\False}{3, \Diamond p}, just = {\TRule{\False}{\Diamond}[11]}, close]
                                    ]
                                ]
                            ]
                        ]
                    ]
                ]
            ]
        ]
    ]
\end{oltableau}

Both branches close, so this is \textbf{valid} in \textbf{S5}. There are many ways to close this one; this is the fastest I found, but it's definitely not the only way.

\newpage

\rulesubheader{$(\Diamond p \wedge \Box q) \rightarrow \Diamond q$ in \textbf{K}}

\begin{oltableau}
    [\sFmla{\False}{1, (\Diamond p \wedge  \Box q) \rightarrow  \Diamond q}, checked, just = \TAss
        [\sFmla{\True}{1, \Diamond p \wedge  \Box q}, checked, just = {\TRule{\False}{\rightarrow}[1]}
            [\sFmla{\False}{1, \Diamond q}, just = {\TRule{\False}{\rightarrow}[1]}
                [\sFmla{\True}{1, \Diamond p}, checked, just = {\TRule{\True}{\wedge}[2]}
                    [\sFmla{\True}{1, \Box q}, just = {\TRule{\True}{\wedge}[2]}
                        [\sFmla{\True}{1.1, p}, just = {\TRule{\True}{\Diamond}[4]}
                            [\sFmla{\False}{1.1, q}, just = {\TRule{\False}{\Diamond}[3]}
                                [\sFmla{\True}{1.1, q}, just = {\TRule{\True}{\Box}[5]}, close]
                            ]
                        ]
                    ]
                ]
            ]
        ]
    ]
\end{oltableau}

The trunk closes, so it is \textbf{valid} in \textbf{K}. What's odd about this one is that even though the right hand side of the conditional is just about $q$, the $\Diamond p$ plays a crucial role in the tableau, by making us introduce $w_1.1$.

\rulesubheader{$\Box (p \wedge \neg p) \vee \Diamond (q \vee \neg q)$ in \textbf{K}}

\begin{oltableau}
    [\sFmla{\False}{1, \Box (p \wedge  \neg p) \lor  \Diamond (q \lor  \neg q)}, checked, just = \TAss
        [\sFmla{\False}{1, \Box (p \wedge  \neg p)}, checked, just = {\TRule{\False}{\lor}[1]}
            [\sFmla{\False}{1, \Diamond (q \lor  \neg q)}, just = {\TRule{\False}{\lor}[1]}
                [\sFmla{\False}{1.1, p \wedge  \neg p}, just = {\TRule{\False}{\Box}[2]}
                    [\sFmla{\False}{1.1, q \lor  \neg q}, checked, just = {\TRule{\False}{\Diamond}[3]}
                        [\sFmla{\False}{1.1, q}, just = {\TRule{\False}{\wedge}[5]}
                            [\sFmla{\False}{1.1, \neg q}, checked, just = {\TRule{\False}{\lor}[5]}
                                [\sFmla{\True}{1.1, q}, just = {\TRule{\False}{\neg}[7]}, close]
                            ]
                        ]
                    ]
                ]
            ]
        ]
    ]
\end{oltableau}

This is also valid, and it's also wrong for the same reason. (Not at all on the test digression. This example shows that \textbf{K} has a very weird property. It is not Halldén-complete. A logician, Sören Halldén, in the 1950s noticed that some logics have the odd property that $X \vee Y$ is a theorem even though (a) neither $X$ nor $Y$ is, and (b) $X$ and $Y$ have no variables in common. We've come to call such logics Halldén-incomplete, and the opposite property Halldén-completeness. In general, the logics that are interesting and get used in philosophical applications are Halldén-complete.)

\rulesubheader{$\Box \neg(p \vee \neg q) \rightarrow \Box p$ in \textbf{K}}

\begin{oltableau}
    [\sFmla{\False}{1, \Box \neg (p \lor  \neg q) \rightarrow  \Box p}, checked, just = \TAss
        [\sFmla{\True}{1, \Box \neg (p \lor  \neg q)}, just = {\TRule{\True}{\rightarrow}[1]}
            [\sFmla{\False}{1, \Box p}, checked, just = {\TRule{\False}{\rightarrow}[1]}
                [\sFmla{\False}{1.1, p}, just = {\TRule{\False}{\Box}[3]}
                    [\sFmla{\True}{1.1, \neg (p \lor  \neg q)}, checked, just = {\TRule{\True}{\Box}[2]}
                        [\sFmla{\False}{1.1, p \lor  \neg q}, checked, just = {\TRule{\True}{\neg}[5]}
                            [\sFmla{\False}{1.1, p}, just = {\TRule{\False}{\lor}[6]}
                                [\sFmla{\False}{1.1, \neg q}, checked, just = {\TRule{\False}{\lor}[6]}
                                    [\sFmla{\True}{1.1, q}, just = {\TRule{\False}{\neg}[8]}]
                                ]
                            ]
                        ]
                    ]
                ]
            ]
        ]
    ]
\end{oltableau}

The tree is open, so this is not a theorem, it is \textbf{invalid} in \textbf{K}. The countermodel has two worlds: $w_1$ and $w_1.1$, with $w_1Rw_{1.1}$. At $w_{1.1}$, $p$ is false and $q$ is true.

\newpage

\rulesubheader{$\Box(p \rightarrow (q \wedge r)) \rightarrow (\Diamond \Diamond p \rightarrow \Diamond (q \vee s))$ in \textbf{K}}

\begin{oltableau}
    [\sFmla{\False}{1, \Box (p \rightarrow  (q \wedge  r)) \rightarrow  (\Diamond \Diamond p \rightarrow  \Diamond (q \lor  s))}, checked, just = \TAss
        [\sFmla{\True}{1, \Box (p \rightarrow  (q  \wedge  r))}, just = {\TRule{\False}{\rightarrow}[1]}
            [\sFmla{\False}{1, \Diamond \Diamond p \rightarrow  \Diamond (q \lor  s)}, checked, just = {\TRule{\False}{\rightarrow}[1]}
                [\sFmla{\True}{1, \Diamond \Diamond p}, checked, just = {\TRule{\False}{\rightarrow}[3]}
                    [\sFmla{\False}{1, \Diamond (q \lor  s)}, just = {\TRule{\False}{\rightarrow}[3]}
                        [\sFmla{\True}{1.1, \Diamond p}, checked, just = {\TRule{\True}{\Diamond}[4]}
                            [\sFmla{\True}{1.1, \Box (p \rightarrow  (q  \wedge  r))}, just = {\TRule{\Box}{\text{4}}[2]}
                                [\sFmla{\False}{1.1, \Diamond (q \lor  s)}, just = {\TRule{\Diamond}{\text{4}}[5]}
                                    [\sFmla{\True}{1.1.1, p}, just = {\TRule{\True}{\Diamond}[6]}
                                        [\sFmla{\True}{1.1.1, p \rightarrow  (q \wedge  r)}, just = {\TRule{\True}{\Box}[7]}
                                            [\sFmla{\False}{1.1.1, p}, just = {\TRule{\True}{\rightarrow}[10]}, close]
                                            [\sFmla{\True}{1.1.1, q \wedge  r}, checked, just = {\TRule{\True}{\rightarrow}[10]}
                                                [\sFmla{\True}{1.1.1, q}, just = {\TRule{\True}{\wedge}[11]}
                                                    [\sFmla{\True}{1.1.1, r}, just = {\TRule{\True}{\wedge}[11]}
                                                        [\sFmla{\False}{1.1.1, q \lor  s}, just = {\TRule{\False}{\Diamond}[8]}
                                                            [\sFmla{\False}{1.1.1, q}, just = {\TRule{\False}{\lor}[14]}, close]
                                                        ]
                                                    ]
                                                ]
                                            ]
                                        ]
                                    ]
                                ]
                            ]
                        ]
                    ]
                ]
            ]
        ]
    ]
\end{oltableau}

All the branches close, so this is valid.

\rulesubheader{$(\Diamond \Box p \wedge \Diamond \Box q) \rightarrow \Diamond(p \wedge q)$ in \textbf{K}}

\begin{oltableau}
    [\sFmla{\False}{1, (p \wedge  \Diamond \Box p) \rightarrow  \Box p}, checked, just = \TAss
        [\sFmla{\True}{1, p \wedge  \Diamond \Box p}, checked, just = {\TRule{\False}{\rightarrow}[1]}
            [\sFmla{\False}{1, \Box p}, checked, just = {\TRule{\False}{\rightarrow}[1]}
                [\sFmla{\True}{1, p}, just = {\TRule{\True}{\wedge}[2]}
                    [\sFmla{\True}{1, \Diamond \Box p}, checked, just = {\TRule{\True}{\wedge}[2]}
                        [\sFmla{\False}{1.1, p}, just = {\TRule{\False}{\Box}[3]}
                            [\sFmla{\True}{1.2, \Box p}, just = {\TRule{\True}{\Diamond}[5]}
                                [\sFmla{\True}{1.2, p}, just = {\TRule{\Box}{\text{T}}[7]}]
                            ]
                        ]
                    ]
                ]
            ]
        ]
    ]
\end{oltableau}

This is open, so it isn't valid. The countermodel has three worlds, $w_1$, $w_{1.1}$ and $w_{1.2}$. $w_1$ can access everything, the other two can only access themselves. And $p$ is true at $w_1$ and $w_{1.2}$, but false at $w_{1.1}$.

\rulesubheader{$ \Box \Diamond \Box \Diamond p \rightarrow \Box \Diamond p$ in \textbf{K}}

\begin{oltableau}
    [\sFmla{\False}{1, \Diamond \Box \Diamond \Box p \rightarrow  \Diamond \Box p}, checked, just = \TAss
        [\sFmla{\True}{1, \Diamond \Box \Diamond \Box p}, checked, just = {\TRule{\False}{\rightarrow}[1]}
            [\sFmla{\False}{1, \Diamond \Box p}, just = {\TRule{\False}{\rightarrow}[1]}
                [\sFmla{\True}{1.1, \Box \Diamond \Box p}, just = {\TRule{\True}{\Diamond}[2]}
                    [\sFmla{\False}{1.1, \Diamond \Box p}, just = {\TRule{\Diamond}{\text{4}}[3]}
                        [\sFmla{\True}{1.1, \Diamond \Box p}, just = {\TRule{\Box}{\text{T}}[4]}, close]
                    ]
                ]
            ]
        ]
    ]
\end{oltableau}

The trunk closes, so this is \textbf{valid}. There are \textit{a lot} of ways to do this one; this is one of the more direct.

\end{document}
